\documentclass[11pt, a4paper, oneside]{article}
\usepackage[ngerman]{babel}
\usepackage{worksheet}

\begin{document}
	\author{L. Bung}
	\title{Die Eulersche Zahl}
	\subject{Mathematik}
	\class{FHR1}
	\maketitle
	
	Investiert man Geld bei einer Bank, bekommt man darauf Zinsen.
	Die Bank hat dabei mehrere Optionen, um die Zinsen auszuzahlen:
	jährlich, halbjährlich, monatlich...
	
	Angenommen, man investiert einen Euro bei 100\% Verzinsung.
	Werden die Zinsen jährlich ausgezahlt, ist der Betrag nach einem Jahr
	$\displaystyle\left(1 + \frac{1}{1}\right)^1$.
	
	Bei monatlicher Verzinsung sind es
	$\displaystyle\left(1 + \frac{1}{12}\right)^{12}$.
	
	Allgemein ist der Betrag bei $n$ Zeitintervallen
	$\displaystyle\left(1 + \frac{1}{n}\right)^n$.
	
	\singletask{Zinszahlung zu verschiedenen Zeitpunkten}
	
	Berechnen Sie jeweils, wie hoch der Endbetrag nach einem Jahr ist:
	
	\begin{table}[h]
		\centering
		\begin{tabular}{|p{4cm}||p{1.5cm}|p{1.5cm}|p{1.5cm}|p{1.5cm}|p{1.5cm}|}
			\hline
			\textbf{Anzahl der Intervalle} & 1 & 2 & 12 & 48 & 365\\
			\hline\hline
			\textbf{Endbetrag} &&&&&\\
			\hline
		\end{tabular}
	\end{table}
	
	\begin{hint}{Die Eulersche Zahl}
		\phantom{x}\vspace{6cm}
	\end{hint}
	
	\bonustask{Graph der e-Funktion}
	
	Zeichnen Sie den Graphen der Funktion $f(x)=e^x$ in ein Koordinatensystem.
	
\end{document}